% Trechos para incorporar ao texto-fonte do TCC. Requer booktabs e amsmath.
% Conferidos com src/baseline_cnn.py e src/train_baseline.py no commit
% 07ee6c3447c48fa05b0c38b24aa46522f0b84fd9.
% O PDF de 22/11/2026 repete as subseções 5.1--5.3 em 6.1--6.3.
% Integrar os trechos uma única vez, após consolidar essas seções.

% Substituir o parágrafo resumido sobre a CNN em Modelos e configurações.
\subsubsection{Arquitetura da CNN própria}
A CNN própria recebe imagens RGB de $224\times224$ pixels e foi treinada
a partir de pesos inicializados sem pré-treinamento. Sua arquitetura possui
quatro blocos, cada um composto, nesta ordem, por uma convolução, uma
ativação ReLU, uma operação de max-pooling e uma camada de dropout espacial
(\texttt{Dropout2d}). As convoluções possuem filtros $3\times3$, passo 1,
preenchimento de 1 pixel e termo de viés. Os blocos utilizam, respectivamente,
32, 64, 128 e 256 filtros. O max-pooling utiliza janela $2\times2$, passo 2 e
sem preenchimento. As probabilidades de dropout espacial são 0,10, 0,15,
0,20 e 0,25, respectivamente.

Após o quarto bloco, aplica-se global average pooling, implementado por
\texttt{AdaptiveAvgPool2d((1,1))}, produzindo 256 características. O tensor
é convertido em um vetor, seguido por dropout com probabilidade 0,50 e uma
camada linear com 38 saídas. A rede contém 398.182 parâmetros, todos
treináveis. As saídas são logits, fornecidos diretamente à
\texttt{CrossEntropyLoss}, sem uma camada softmax explícita no modelo.
A classe predita é determinada pelo maior logit. Durante a avaliação, o
modelo opera em modo de inferência, com dropout desativado.

\begin{table*}[t]
\centering
\caption{Estrutura da CNN própria. As dimensões representam canais, altura
e largura, sem a dimensão do lote.}
\label{tab:cnn-propria-detalhada}
\small
\begin{tabular}{lrrr}
\toprule
Etapa & Saída & Dropout & Parâmetros \\
\midrule
Entrada & $3\times224\times224$ & -- & 0 \\
Bloco 1 & $32\times112\times112$ & 0,10 & 896 \\
Bloco 2 & $64\times56\times56$ & 0,15 & 18.496 \\
Bloco 3 & $128\times28\times28$ & 0,20 & 73.856 \\
Bloco 4 & $256\times14\times14$ & 0,25 & 295.168 \\
Global average pooling & $256\times1\times1$ & -- & 0 \\
Flatten e dropout & 256 & 0,50 & 0 \\
Camada linear & 38 & -- & 9.766 \\
\midrule
Total & -- & -- & 398.182 \\
\bottomrule
\end{tabular}
\end{table*}

% Acrescentar aos detalhes de treinamento existentes; evitar repeti-los.
\subsubsection{Detalhes do treinamento da CNN própria}
O treinamento da CNN própria utilizou Adam com taxa de aprendizado fixa
de $10^{-3}$, função de perda \texttt{CrossEntropyLoss}, lotes de 32
imagens e limite de 30 épocas. Não foi utilizado agendador de taxa de
aprendizado. As imagens foram normalizadas com médias RGB
$(0{,}485,0{,}456,0{,}406)$ e desvios padrão
$(0{,}229,0{,}224,0{,}225)$. No conjunto de treinamento, aplicaram-se
rotações aleatórias no intervalo de $-15^\circ$ a $15^\circ$, espelhamento
horizontal com probabilidade 0,5 e variações de brilho e contraste com
fatores entre 0,9 e 1,1. A ordem das imagens foi embaralhada no treinamento.
Na validação e no teste, foram utilizados apenas redimensionamento,
conversão para tensor e normalização.

A parada antecipada acompanhou a perda de validação, com paciência de
cinco épocas e melhoria mínima de $10^{-4}$. A gravação do checkpoint
utilizou um critério separado: qualquer redução estrita na perda de
validação. O checkpoint com menor perda foi restaurado antes da avaliação.
Nas execuções com sementes 42, 123 e 456, a CNN própria completou as 30
épocas, e os checkpoints selecionados corresponderam às épocas 29, 29 e
30, respectivamente. A divisão dos dados permaneceu fixa em 38.008 imagens
de treinamento, 8.172 de validação e 8.125 de teste, com separação por
\texttt{leaf\_id}. Os parâmetros de treinamento foram definidos sem utilizar
o conjunto de teste para seleção.

% Acrescentar ao protocolo de avaliação.
\subsubsection{Comparação estatística}
Foi aplicado o teste de postos sinalizados de Wilcoxon pareado, bilateral
\cite{wilcoxon1945},
às métricas de acurácia e F1-score macro do conjunto de teste. Para cada
par de arquiteturas, os resultados foram pareados pelas sementes 42, 123
e 456, totalizando três pares por comparação. A unidade de análise foi a
execução de treinamento; as imagens, classes e épocas não foram tratadas
como repetições independentes. As seis combinações de arquiteturas e as
duas métricas constituíram uma família de 12 testes, à qual se aplicou a
correção de Holm \cite{holm1979}, com nível de significância de 5\%.

O cálculo utilizou \texttt{scipy.stats.wilcoxon}, com p-valores exatos.
Na comparação de acurácia entre DenseNet201 e EfficientNet-B0, que apresentou
empate nos valores absolutos das diferenças, utilizou-se a enumeração
exaustiva das permutações de sinais. As diferenças foram arredondadas a 12 casas decimais
para evitar artefatos numéricos na ordenação dos postos. A interpretação
do teste pressupõe diferenças pareadas independentes e uma distribuição
simétrica sob a hipótese nula; apenas três pares não permitem avaliar
adequadamente essa simetria. A análise é condicionada à divisão fixa dos
dados e não estima a variabilidade entre diferentes conjuntos de teste.

% Acrescentar em Resultados e discussão, após os resultados globais.
\subsection{Resultados da comparação estatística}
A Tabela~\ref{tab:wilcoxon-arquiteturas} apresenta os p-valores brutos do
teste de Wilcoxon. Após a correção de Holm sobre os 12 testes, todos os
p-valores ajustados foram iguais a 1,00. Não houve rejeição da hipótese
nula ao nível de 5\% em nenhuma das comparações.

\begin{table*}[t]
\centering
\caption{P-valores brutos do Wilcoxon bilateral, com três pares por
comparação. Todos os p-valores ajustados por Holm foram 1,00.}
\label{tab:wilcoxon-arquiteturas}
\small
\begin{tabular}{lrr}
\toprule
Par de arquiteturas & Acurácia & F1 macro \\
\midrule
CNN própria e ResNet50 & 0,25 & 0,25 \\
CNN própria e DenseNet201 & 0,25 & 0,25 \\
CNN própria e EfficientNet-B0 & 0,25 & 0,25 \\
ResNet50 e DenseNet201 & 1,00 & 0,75 \\
ResNet50 e EfficientNet-B0 & 0,25 & 0,50 \\
DenseNet201 e EfficientNet-B0 & 0,25 & 0,25 \\
\bottomrule
\end{tabular}
\end{table*}

Esse resultado deve ser interpretado considerando o número reduzido de
execuções. Com três diferenças não nulas, existem apenas $2^3=8$
configurações de sinais, e o menor p-valor bilateral exato possível é
$2/8=0{,}25$, mesmo quando uma arquitetura supera a outra nas três
sementes. Assim, as diferenças de desempenho observadas permanecem
descritivas: a ausência de significância não demonstra igualdade ou
equivalência entre as arquiteturas. A EfficientNet-B0 apresentou as maiores
médias observadas de acurácia e F1 macro, sem que este teste permita
afirmar superioridade estatística.

% Substituir, na conclusão, a afirmação de que não foram realizados testes.
% Sugestão de frase:
% A comparação utilizou uma divisão fixa dos dados e três sementes de
% treinamento. O teste bilateral de Wilcoxon pareado, com correção de Holm,
% não identificou diferenças estatisticamente significativas a 5\%, resultado
% limitado pela resolução do teste com apenas três pares e que não demonstra
% equivalência entre as arquiteturas.
